\section{The positive crossing kernel}\label{sec:crossing}

Let \(F_\mu,F_\nu\) be the CDFs of \eqref{eq:mu}--\eqref{eq:nu}, and put
\[
 \Delta(y)=F_\mu(y)-F_\nu(y).
\]

\begin{lemma}[Strict stochastic ordering]\label{lem:crossing}
One has \(\Delta(y)>0\) for \(p<y<w\), while
\(\Delta(p)=\Delta(w)=0\).
\end{lemma}

\begin{proof}
On \([z,w]\), \(F_\mu=1\), so \(\Delta=1-F_\nu>0\) before \(w\). On
\((p,z)\), the ratio of the two positive densities is
\begin{equation}\label{eq:density-ratio}
 R_0(y)=\frac{\dd\mu/\dd y}{\dd\nu/\dd y}
 =\frac{\Lambda(z-y)}{L\delta(y-p)}.
\end{equation}
Writing \(C=\Lambda/(L\delta)>0\), direct differentiation gives
\begin{equation}\label{eq:density-ratio-derivative}
 R_0'(y)=-\frac{C(z-p)}{(y-p)^2}<0.
\end{equation}
Thus \(R_0\) decreases strictly from \(+\infty\) to \(0\). Hence
\(\Delta'=\mu'-\nu'\) changes sign exactly once, from positive to negative.
Since \(\Delta(p)=0\) and
\[
 \Delta(z)=1-F_\nu(z)=\nu((z,w])>0,
\]
the function cannot return to zero on \((p,z]\). The right interval was
already handled.
\end{proof}

Integration by parts for CDFs, using \(A_0'=D\) and the common total mass,
turns Lemma~\ref{lem:transport} into
\begin{equation}\label{eq:crossing-integral}
 K=\int_p^w \Delta(y)D(y)\dd y.
\end{equation}
Combining \eqref{eq:K-def}, \eqref{eq:C4-curvature}, and
\eqref{eq:crossing-integral} gives the central identity
\begin{equation}\label{eq:central-identity}
 \mathcal N=\delta\Lambda\int_p^w
 \Delta(y)\frac{\Cfour(y)}{Q(y)^6\kappa(y)^2}\dd y>0.
\end{equation}
Every factor in the interior is strictly positive.
